% SEGMENT OLP-0609-S01
% Part: methods
% Chapter: proofs
% Section: proof-by-contradiction

\documentclass[../../../include/open-logic-section]{subfiles}

\begin{document}

\olfileid{mth}{prf}{con}

\olsection{முரண்பாட்டால் நிறுவுதல்}

முதலில், முரண்பாட்டால் நிறுவுதல் என்பது எதிர்மறைக் கூற்றுகளை
நிறுவப் பயன்படும் உய்த்தறிதல் வடிவம். ஒரு கூற்று $p$
\emph{பொய்} என்று, அதாவது $\lnot p$ என்று, காட்ட
விரும்புகிறீர்கள் என்க. பயனுள்ள உத்தி: (a) $p$ மெய்
என எடுத்துக்கொண்டு, (b) இந்த எடுகோள் பொய்யென்று
ஏற்கெனவே தெரிந்த ஒன்றுக்கு இட்டுச் செல்கிறது என்று
காட்டுவது. அந்தப் பொய்யென்று தெரிந்தது $p$ உடனே
முரண்படும் முடிவாகவோ, ஆய்விலுள்ள மொத்தக் கூற்றின்
மற்றோர் எடுகோளுடன் முரண்படும் முடிவாகவோ இருக்கலாம்.
எடுத்துக்காட்டாக, ``$q$ எனில் $\lnot p$'' என்பதை
நிறுவ, $q$ மெய் என்று எடுத்துக்கொண்டு அதிலிருந்து
$\lnot p$ ஐ நிறுவ வேண்டும். $\lnot p$ ஐ
முரண்பாட்டால் நிறுவினால், $q$ உடன் $p$ ஐயும்
எடுத்துக்கொள்ள வேண்டும். $p$ இலிருந்து $\lnot q$
ஐ நிறுவ முடிந்தால், $p$ என்ற எடுகோள் மற்றோர்
எடுகோளான $q$ உடன் முரண்படுகிறது: $q$ உம்
$\lnot q$ உம் ஒரே நேரத்தில் மெய்யாக முடியாது.
முரண்பாட்டை அடைவதற்கான நிறுவலில் பிற உய்த்தறிதல்
வடிவங்களையும் வரையறைகளை விரித்துரைப்பதையும் பயன்படுத்த
வேண்டியிருக்கும். ஓர் எடுத்துக்காட்டைப் பார்ப்போம்.

\begin{prop}
  $A \subseteq B$ மற்றும் $B = \emptyset$ எனில்,
  $A$ க்கு !!{element}s இல்லை.
\end{prop}

% SEGMENT OLP-0609-S02
\begin{proof}
  $A \subseteq B$ மற்றும் $B = \emptyset$ என்க.
  $A$ க்கு !!{element}s இல்லை என்று காட்ட வேண்டும்.
  \begin{quote}
    இது நிபந்தனைக் கூற்று என்பதால் அதன் முற்பகுதியை
    எடுத்துக்கொண்டு பிற்பகுதியை நிறுவ வேண்டும்.
    பிற்பகுதி: $A$ க்கு !!{element}s இல்லை.
    இதை இன்னும் வெளிப்படையாகக் கூறலாம்:
    ஓர் $x \in A$ உள்ளது என்பது மெய்யல்ல.
  \end{quote}
  $A$ க்கு !!{element}s இல்லை எனில், எனில் மட்டுமே,
  $x \in A$ ஆகும் ஓர் $x$ உள்ளது என்பது மெய்யல்ல.
  \begin{quote}
    நிறுவ வேண்டியது உண்மையில் $\lnot p$ என்ற
    எதிர்மறைக் கூற்று என்பதை அறிந்தோம்: ஓர் $x \in A$
    உள்ளது என்பது மெய்யல்ல. முரண்பாட்டால் நிறுவ,
    அதற்குரிய நேர்மறைக் கூற்றான $p$ ஐ, அதாவது
    ஓர் $x \in A$ உள்ளது என்று, எடுத்துக்கொண்டு
    அதிலிருந்து முரண்பாட்டை நிறுவ வேண்டும்.
    ``முரண்பாட்டிற்காக எடுத்துக்கொள்க'' அல்லது
    ``அப்படியல்ல என்க'' என்று எழுதி, பின் $p$
    என்ற எடுகோளைச் சொல்வதன் மூலம் இம்முறையைக்
    குறிக்கிறோம்.
  \end{quote}
  அப்படியல்ல என்க: ஓர் $x \in A$ உள்ளது.
  \begin{quote}
    இப்போது முரண்பாட்டைப் பெற இந்தப் புதிய எடுகோளைப்
    பயன்படுத்துவோம். மேலும் இரு எடுகோள்கள் உள்ளன:
    $A \subseteq B$ மற்றும் $B = \emptyset$.
    முதலாவது $x \in B$ என்பதைக் கொடுக்கிறது:
  \end{quote}
  $A \subseteq B$ என்பதால் $x \in B$.
  \begin{quote}
    ஆனால் $B = \emptyset$ என்பதால் $B$ இன்
    ஒவ்வோர் !!{element} உம், எடுத்துக்காட்டாக $x$,
    $\emptyset$ இன் !!a{element} ஆகவும் இருக்க வேண்டும்.
  \end{quote}
  $B = \emptyset$ என்பதால் $x \in \emptyset$.
  இது முரண்பாடு; வரையறையின்படி $\emptyset$ க்கு
  !!{element}s இல்லை.
  \begin{quote}
    நிறுவல் இப்போதே முடிந்தது: எடுத்த எடுகோள்களிலிருந்து
    தேவைப்பட்ட முரண்பாட்டை அடைந்தோம். எனவே அவையனைத்தும்
    ஒரே நேரத்தில் மெய்யாக முடியாது. முதல் இரு எடுகோள்களான
    $A \subseteq B$, $B = \emptyset$ ஆகியவை
    கேள்விக்குள்ளாக்கப்படாததால், கடைசியாக எடுத்த
    எடுகோளான ஓர் $x \in A$ உள்ளது என்பதே பொய்யாக
    வேண்டும். முழுமையாகக் கூற விரும்பினால் இதை
    வெளிப்படையாக எழுதலாம்.
  \end{quote}
  ஆகவே ஓர் $x \in A$ உள்ளது என்ற எடுகோள்
  பொய்யாக வேண்டும்; இதனால் முரண்பாட்டால் நிறுவியபடி
  $A$ க்கு !!{element}s இல்லை.
\end{proof}

% SEGMENT OLP-0609-S03
ஒவ்வொரு நேர்மறைக் கூற்றும் ஓர் எதிர்மறைக் கூற்றுக்குச்
சமமானது: $p$ எனில், எனில் மட்டுமே, $\lnot\lnot p$.
ஆகவே முரண்பாட்டால் நிறுவுதலை நேர்மறைக் கூற்றுகளுக்கும்
``மறைமுகமாகப்'' பயன்படுத்தலாம். $p$ ஐ நிறுவ,
அதை $\lnot\lnot p$ என்ற எதிர்மறைக் கூற்றாகக்
கருதுக. $\lnot p$ இலிருந்து முரண்பாட்டை நிறுவ
முடிந்தால், முரண்பாட்டால் $\lnot\lnot p$ ஐயும்,
அதனால் $p$ ஐயும் நிறுவியிருக்கிறோம்.

கடைசி எடுத்துக்காட்டில் $A$ க்கு !!{element}s இல்லை
என்ற எதிர்மறைக் கூற்றை நிறுவ விரும்பினோம். எனவே
முரண்பாட்டிற்காக ஓர் $x \in A$ உள்ளது என்ற
நேர்மறைக் கூற்றை எடுத்துக்கொண்டோம். அதிலிருந்து
தற்காலிகமான $x \in A$ பற்றி தொடர்ந்து சிந்தித்து,
இறுதியில் $x \in \emptyset$ என்ற முரண்பாட்டை
அடைந்தோம்.

நேர்மறைக் கூற்றை மறைமுகமாக நிறுவும்போது, முரண்பாட்டிற்காக
எடுக்கும் கூற்று எதிர்மறையாக இருக்கும். ஆனால் பல
நேரங்களில் நேர்மறைக் கூற்றை எதிர்மறையாகவும்,
எதிர்மறைக் கூற்றை நேர்மறையாகவும் எளிதில் மாற்றிக்
கூறலாம். முந்தைய நிறுவலில் ``$A$ க்கு
!!{element}s இல்லை'' என்ற எதிர்மறைப் பிற்பகுதிக்குப்
பதிலாக ``$A = \emptyset$'' என்று நிறுவியிருந்தாலும்
அடிப்படையில் அதே நிறுவல்தான் கிடைத்திருக்கும்.
($=$ இன் வரையறையின்படி, ``$A = \emptyset$'' என்பது
ஒரு பொதுக் கூற்று: ``$A$ இன் ஒவ்வோர் !!{element}
உம் $\emptyset$ இன் !!{element} ஆகவும், மறுதிசையிலும்
இருக்கும்'' என்று விரியும்.) ஆனால் அது ``ஓர்
$x \in A$ உள்ளது என்பது மெய்யல்ல'' என்ற
எதிர்மறைக் கூற்றுக்குச் சமம் என்பது எளிதில் தெரியும்.

சில நேரங்களில் $p$ ஐ நேரடியாக நிறுவுவதைவிட
$\lnot p$ ஐ எடுகோளாகக் கொண்டு வேலை செய்வது
எளிது. நேரடி நிறுவல் அதே அளவுக்கு எளிதாகவோ
மேலும் எளிதாகவோ இருந்தாலும் (அடுத்த எடுத்துக்காட்டுகள்
போல), சிலர் மறைமுகமாகச் செல்ல விரும்புவார்கள்.
இரட்டை மறுப்பு குழப்பம் தந்தால், ஒரு கூற்றை
முரண்பாட்டால் நிறுவுவது அதன் \emph{எதிர்க்} கூற்றிலிருந்து
முரண்பாட்டை நிறுவுவது என்று கருதுக. ஆகவே $\lnot p$
ஐ முரண்பாட்டால் நிறுவுவது $p$ என்ற எடுகோளிலிருந்து
முரண்பாட்டை நிறுவுவதாகும்; $p$ ஐ முரண்பாட்டால்
நிறுவுவது $\lnot p$ என்ற எடுகோளிலிருந்து
முரண்பாட்டை நிறுவுவதாகும்.

% SEGMENT OLP-0609-S04
\begin{prop}
$A \subseteq A \cup B$.
\end{prop}

\begin{proof}
  $A \subseteq A \cup B$ என்று காட்ட வேண்டும்.
  \begin{quote}
    முதல் பார்வையில் இது நேர்மறைக் கூற்று: ஒவ்வொரு
    $x \in A$ உம் $A \cup B$ இல் உள்ளது.
    இதன் மறுப்பு: ஏதோ ஓர் $x \in A$ ஆனது
    $\notin A \cup B$. இந்த மறுத்த கூற்றை
    எடுத்துக்கொண்டு அது முரண்பாட்டிற்கு இட்டுச்
    செல்கிறது என்று காட்டுவதால் மறைமுகமாக நிறுவலாம்.
  \end{quote}
  அப்படியல்ல என்க; அதாவது $A \nsubseteq A \cup B$.
  \begin{quote}
    $A \subseteq A \cup B$ என்பதன் வரையறை:
    ஒவ்வொரு $x \in A$ உம் $\in A \cup B$.
    $A \nsubseteq A \cup B$ என்பதன் பொருளைப்
    புரிந்துகொள்ள, விரித்த வரையறையில் ஓர் எளிய
    தருக்க மாற்றம் தேவை: ஒவ்வொரு $x \in A$
    உம் $\in A \cup B$ என்பது பொய் எனில்,
    எனில் மட்டுமே, ஏதோ \emph{ஓர்} $x \in A$
    % SOURCE CORRECTION TA-MTH-004: தொடர்பில்லாத C ஐ இலக்குக் கணம் A ∪ B ஆகத் திருத்துக.
    ஆனது $\notin A \cup B$. (இங்கே கவனமாக
    இருக்க வேண்டும்: ``எல்லா $A$ வகைப் பொருள்களும் $B$ வகைப் பொருள்களே'' போன்ற
    கூற்றின் மறுப்பைத் தவறாகப் புரிந்துகொள்வதால்
    மாணவர்களின் பல முரண்பாட்டு நிறுவல்கள் தோற்கின்றன.)
    வேறுவிதமாக, $A \nsubseteq A \cup B$ எனில்,
    எனில் மட்டுமே, $x \in A$ மற்றும்
    $x \notin A \cup B$ ஆகும் ஓர் $x$ உள்ளது.
    இனி நிறுவல் எளிது.
  \end{quote}
  ஆகவே $x \notin A \cup B$ ஆகும் ஓர் $x \in A$
  உள்ளது. $\cup$ இன் வரையறையின்படி
  $x \in A \cup B$ எனில், எனில் மட்டுமே,
  $x \in A$ அல்லது $x \in B$. $x \in A$
  என்பதால் $x \in A \cup B$. இது
  $x \notin A \cup B$ என்ற எடுகோளுடன் முரண்படுகிறது.
\end{proof}

\begin{prob}
$A \cap B \subseteq A$ என்பதை \emph{மறைமுகமாக}
நிறுவுக.
\end{prob}

% SEGMENT OLP-0609-S05
\begin{prop}
$A \subseteq B$ மற்றும் $B \subseteq C$ எனில்
$A \subseteq C$.
\end{prop}

\begin{proof}
  $A \subseteq B$ மற்றும் $B \subseteq C$
  என்க. $A \subseteq C$ என்று காட்ட வேண்டும்.
  \begin{quote}
    மறைமுகமாகச் செல்வோம்: நிறுவ விரும்புவதன்
    மறுப்பை எடுகோளாக எடுத்துக்கொள்வோம்.
  \end{quote}
  அப்படியல்ல என்க; அதாவது $A \nsubseteq C$.
  \begin{quote}
    முன்புபோல, $A \nsubseteq C$ எனில்,
    எனில் மட்டுமே, ஒவ்வொரு $x \in A$ உம்
    $\in C$ என்பது பொய்; அதாவது ஏதோ ஓர்
    $x \in A$ ஆனது $\notin C$ என்று
    சிந்திக்கிறோம். பயிற்சி வந்தபின் இத்தகைய
    மறுப்புகளை விரிக்க அதிகம் யோசிக்க வேண்டியிருக்காது.
  \end{quote}
  வேறுவிதமாகச் சொன்னால், $x \in A$ மற்றும்
  $x \notin C$ ஆகும் ஓர் $x$ உள்ளது.
  \begin{quote}
    இதிலிருந்து இப்போது முரண்பாட்டை அடையலாம்.
    அதற்கு மற்ற இரு எடுகோள்களையும் பயன்படுத்த வேண்டும்.
  \end{quote}
  $A \subseteq B$ என்பதால் $x \in B$.
  $B \subseteq C$ என்பதால் $x \in C$.
  ஆனால் இது $x \notin C$ உடன் முரண்படுகிறது.
\end{proof}

% SEGMENT OLP-0609-S06
\begin{prop}
$A \cup B = A \cap B$ எனில் $A = B$.
\end{prop}

\begin{proof}
  $A \cup B = A \cap B$ என்க. $A = B$
  என்று காட்ட வேண்டும்.
  \begin{quote}
    தொடக்கம் இப்போது வழக்கமானதே:
  \end{quote}
  முரண்பாட்டிற்காக $A \neq B$ என எடுத்துக்கொள்க.
  \begin{quote}
    முரண்பாட்டு நிறுவலுக்கான எடுகோள் $A \neq B$.
    $A = B$ எனில், எனில் மட்டுமே, $A \subseteq B$
    மற்றும் $B \subseteq A$ என்பதால், $A \neq B$
    எனில், எனில் மட்டுமே, $A \nsubseteq B$
    \emph{அல்லது} $B \nsubseteq A$. (மறுப்புகளைக்
    கையாளும்போது எவ்வளவு கவனம் தேவை என்று பாருங்கள்!{})
    இந்தப் பிரிப்புக் கூற்றிலிருந்து முரண்பாட்டை நிறுவ,
    இரு நிகழ்வுகளாகப் பிரித்து, ஒவ்வொன்றிலும்
    முரண்பாடு பின்தொடர்கிறது என்று காட்டுகிறோம்.
  \end{quote}
  $A \neq B$ எனில், எனில் மட்டுமே,
  $A \nsubseteq B$ அல்லது $B \nsubseteq A$.
  நிகழ்வுகளைப் பிரிக்கிறோம்.
  \begin{quote}
    முதல் நிகழ்வில் $A \nsubseteq B$ என்க:
    ஏதோ ஓர் $x$ க்கு $x \in A$ ஆனால்
    $\notin B$. $A \cap B$ என்பது $A$, $B$
    இரண்டுக்கும் பொதுவான !!{element}s கொண்ட கணம்.
    ஆகவே ஒரு பொருள் அவற்றில் ஒன்றில் இல்லாவிட்டால்,
    வெட்டிலும் இல்லை. $A \cup B$ என்பது
    $A$ உடன் $B$ ஐ இணைத்த கணம்; ஏதேனும்
    ஒன்றில் உள்ள பொருள் ஒன்றிப்பிலும் இருக்கும்.
    எனவே $x \in A \cup B$ ஆனால்
    $x \notin A \cap B$. இதிலிருந்து
    $A \cap B \neq A \cup B$.
  \end{quote}

  நிகழ்வு 1: $A \nsubseteq B$. அப்போது
  ஏதோ ஓர் $x$ க்கு $x \in A$ ஆனால்
  $x \notin B$. $x \notin B$ என்பதால்
  $x \notin A \cap B$. $x \in A$ என்பதால்
  $x \in A \cup B$. ஆகவே
  $A \cap B \neq A \cup B$; இது
  $A \cap B = A \cup B$ என்ற எடுகோளுடன்
  முரண்படுகிறது.

  நிகழ்வு 2: $B \nsubseteq A$. அப்போது
  ஏதோ ஓர் $y$ க்கு $y \in B$ ஆனால்
  $y \notin A$. முன்புபோல
  $y \in A \cup B$ ஆனால் $y \notin A \cap B$.
  ஆகவே $A \cap B \neq A \cup B$;
  இது மீண்டும் $A \cap B = A \cup B$
  உடன் முரண்படுகிறது.
\end{proof}

\end{document}
